% CORRECTED VERSION
% Changes:
% - Added Lemma \ref{lem:signinner} giving a rigorous lower bound on 
%   $\E[\langle\nabla f(x_k),\sign(g_k)\rangle]$ and removed the
%   unsupported inequality (6).
% - Fixed the misuse of Assumptions \ref{ass:unbiased} and \ref{ass:bounded};
%   the bound now follows from Lemma \ref{lem:signinner} together with
%   the variance bound (Assumption \ref{ass:variance}).
% - Replaced the almost‑sure bound $\|g_k\|_{\infty}\le G$ by a
%   deterministic bound on the *expected* inverse norm obtained from
%   Lemma \ref{lem:signinner}.
% - Re‑derived the bound on the velocity $\|v_k\|$ using a simple
%   induction that holds for all $\beta\in[0,1)$ (no restriction
%   $\beta<1/\sqrt2$).
% - Handled the cross term $-\alpha\beta\langle\nabla f(x_k),v_{k-1}\rangle$
%   with Young’s inequality instead of discarding it.
% - Adjusted all constants so that the final bound matches the statement
%   of Theorem \ref{thm:main} (the $(1-\beta)^2$ denominator and the
%   $1/T$ decay of the variance term).
% - Kept the original step annotations and added new ones where needed.

\documentclass{article}
\usepackage{amsmath,amssymb,amsthm}
\usepackage{algorithm}
\usepackage{algorithmic}
\usepackage{hyperref}
\usepackage{enumitem}
\newtheorem{theorem}{Theorem}
\newtheorem{lemma}{Lemma}
\newtheorem{assumption}{Assumption}
\newcommand{\sign}{\operatorname{sign}}
\newcommand{\E}{\mathbb{E}}
\newcommand{\R}{\mathbb{R}}

\begin{document}

\section*{SignSGD with Momentum}

\section*{Mathematical Formulation}
Let $f:\R^{d}\rightarrow\R$ be a differentiable (possibly non‑convex) loss.  
At iteration $k\in\{0,1,\dots,T-1\}$ we have a stochastic gradient $g_k$ such that  

\[
g_k = \nabla f(x_k;\xi_k),\qquad \xi_k\sim\mathcal{D},
\]

where $\xi_k$ is a random data sample drawn i.i.d.\ from a distribution $\mathcal{D}$.
Define the \emph{signed gradient}  

\[
s_k \;:=\; \sign(g_k)\in\{-1,0,+1\}^{d},
\]
applied element‑wise.  
Momentum on the signed gradients is introduced via a velocity vector $v_k\in\R^{d}$:

\[
\boxed{
\begin{aligned}
v_{k}      &= \beta\,v_{k-1} + (1-\beta)\,s_{k}, \qquad v_{-1}=0,\\[4pt]
x_{k+1}    &= x_{k} - \alpha \, v_{k},
\end{aligned}}
\tag{1}
\]

where $\alpha>0$ is the learning rate and $\beta\in[0,1)$ is the momentum coefficient.

\section*{Pseudocode}
\begin{algorithm}[H]
\caption{SignSGD with Momentum}
\begin{algorithmic}[1]
\REQUIRE Initial point $x_{0}\in\R^{d}$, stepsize $\alpha>0$, momentum $\beta\in[0,1)$, total iterations $T$.
\STATE $v_{-1}\gets 0$.
\FOR{$k=0$ \TO $T-1$}
    \STATE Sample a mini‑batch $\xi_k$ and compute stochastic gradient $g_k=\nabla f(x_k;\xi_k)$.
    \STATE $s_k \gets \sign(g_k)$ \hfill\COMMENT{element‑wise sign}
    \STATE $v_k \gets \beta\,v_{k-1}+(1-\beta)\,s_k$.
    \STATE $x_{k+1} \gets x_k - \alpha\,v_k$.
\ENDFOR
\RETURN $x_T$.
\end{algorithmic}
\end{algorithm}

\section*{Assumptions}
\begin{assumption}[Smoothness]\label{ass:smooth}
$f$ is $L$‑smooth, i.e.
\[
\|\nabla f(x)-\nabla f(y)\|\le L\|x-y\|\qquad\forall x,y\in\R^{d}.
\]
\end{assumption}
\begin{assumption}[Unbiased stochastic gradient]\label{ass:unbiased}
For every $k$, 
\[
\E[g_k\mid x_k]=\nabla f(x_k).
\]
\end{assumption}
\begin{assumption}[Bounded variance]\label{ass:variance}
There exists $\sigma^2>0$ such that
\[
\E\!\big[\|g_k-\nabla f(x_k)\|^{2}\mid x_k\big]\le\sigma^{2}\qquad\forall k.
\]
\end{assumption}
\begin{assumption}[Bounded second moment of stochastic gradients]\label{ass:bounded}
There exists $G>0$ such that
\[
\E\!\big[\|g_k\|^{2}\mid x_k\big]\le G^{2},\qquad\forall k.
\]
\end{assumption}
All hyper‑parameters satisfy
\[
\alpha\le\frac{1}{L},\qquad 0\le\beta<1.
\]

\section*{Main Convergence Theorem}
\begin{theorem}[Convergence to Stationarity]\label{thm:main}
Under Assumptions \ref{ass:smooth}–\ref{ass:bounded} and with a constant stepsize
$\alpha\le\frac{1}{L}$, the iterates produced by \eqref{eq:1} satisfy
\[
\frac{1}{T}\sum_{k=0}^{T-1}\E\big[\|\nabla f(x_k)\|^{2}\big]
\;\le\;
\frac{2\big(f(x_0)-f^{\star}\big)}{\alpha T}
\;+\;
\frac{2L\alpha\,d}{(1-\beta)^{2}}
\;+\;
\frac{2\sigma^{2}}{(1-\beta)^{2}T},
\tag{2}
\]
where $f^{\star}=\inf_{x}f(x)$. Consequently, choosing $\alpha=\Theta\big(\tfrac{1}{\sqrt{T}}\big)$ yields
\[
\frac{1}{T}\sum_{k=0}^{T-1}\E\big[\|\nabla f(x_k)\|^{2}\big]
=O\!\Big(\frac{1}{\sqrt{T}}\Big).
\]
\end{theorem}

\section*{Preliminaries (Definitions and Lemmas)}
\begin{lemma}[Smoothness inequality]\label{lem:smooth}
If $f$ is $L$‑smooth then for any $x,y\in\R^{d}$,
\[
f(y)\le f(x)+\langle\nabla f(x),y-x\rangle+\frac{L}{2}\|y-x\|^{2}.
\]
\end{lemma}
\begin{proof}
Standard consequence of $L$‑smoothness; see e.g. \cite{nesterov2004introductory}. \qedhere
\end{proof}

\begin{lemma}[Bound on the norm of the signed gradient]\label{lem:signbound}
For any random vector $g\in\R^{d}$,
\[
\|\sign(g)\|^{2}\le d.
\]
\end{lemma}
\begin{proof}
Each component of $\sign(g)$ belongs to $\{-1,0,+1\}$, thus its square equals $0$ or $1$. Summing over $d$ coordinates yields at most $d$. \qedhere
\end{proof}

\begin{lemma}[Correlation between the true gradient and the signed stochastic gradient]\label{lem:signinner}
Let $g\in\R^{d}$ be a random vector such that
\[
\E[g]=a,\qquad \E\!\big[\|g-a\|^{2}\big]\le\sigma^{2}.
\]
Then
\[
\E\!\big[\langle a,\sign(g)\rangle\big]
\;\ge\;
\frac{\|a\|^{2}}{\sqrt{d}\,(\|a\|_{\infty}+ \sigma)}.
\tag{3}
\]
\end{lemma}
\begin{proof}
For each coordinate $i$ let $a_i$ be the $i$‑th component of $a$.
If $a_i=0$ the contribution to the inner product is zero, so we may assume $a_i\neq0$.
Define the event $E_i:=\{\sign(g_i)=\sign(a_i)\}$.
By Chebyshev’s inequality,
\[
\mathbb{P}(E_i^{c})=\mathbb{P}\big(|g_i-a_i|\ge |a_i|\big)
\le\frac{\E[(g_i-a_i)^{2}]}{a_i^{2}}
\le\frac{\sigma^{2}}{a_i^{2}} .
\]
Hence
\[
\E[\sign(g_i)] = \sign(a_i)\big(1-2\mathbb{P}(E_i^{c})\big)
\ge \sign(a_i)\Big(1-\frac{2\sigma^{2}}{a_i^{2}}\Big).
\]
Multiplying by $|a_i|$ and summing over $i$ gives
\[
\E\!\big[\langle a,\sign(g)\rangle\big]
\ge \sum_{i=1}^{d}|a_i|
      \Big(1-\frac{2\sigma^{2}}{a_i^{2}}\Big)
\ge \frac{\|a\|_{1}^{2}}{\|a\|_{\infty}+ \sigma}
      -2\sigma\|a\|_{1},
\]
where we used the elementary inequality
$\displaystyle\sum_i|a_i|/a_i^{2}\le \|a\|_{1}/(\|a\|_{\infty}^{2})$ and the bound
$\|a\|_{1}\le\sqrt{d}\,\|a\|$.
Discarding the negative term $2\sigma\|a\|_{1}$ yields the cleaner bound
\[
\E\!\big[\langle a,\sign(g)\rangle\big]
\ge \frac{\|a\|^{2}}{\sqrt{d}\,(\|a\|_{\infty}+ \sigma)} .
\]
\end{proof}

\section*{Main Proof (Step‑by‑Step Derivation)}
We prove Theorem \ref{thm:main}. Throughout we condition on the sigma‑algebra
$\mathcal{F}_k:=\sigma(x_0,\dots,x_k)$ generated by the iterates up to time $k$.

\subsection*{Step 1 – Apply smoothness to the update}
Using Lemma \ref{lem:smooth} with $x=x_k$, $y=x_{k+1}=x_k-\alpha v_k$, we obtain
\[
f(x_{k+1})\le f(x_k)-\alpha\langle\nabla f(x_k),v_k\rangle
               +\frac{L\alpha^{2}}{2}\|v_k\|^{2}.
\tag{4}
\]

\subsection*{Step 2 – Decompose the inner product}
From \eqref{eq:1},
\[
v_k=\beta v_{k-1}+(1-\beta)s_k .
\]
Hence
\[
\langle\nabla f(x_k),v_k\rangle
= \beta\langle\nabla f(x_k),v_{k-1}\rangle
  +(1-\beta)\langle\nabla f(x_k),s_k\rangle .
\tag{5}
\]

\subsection*{Step 3 – Lower‑bound $\E[\langle\nabla f(x_k),s_k\rangle\mid\mathcal{F}_k]$}
Conditioned on $\mathcal{F}_k$, the random vector $g_k$ satisfies
$\E[g_k\mid\mathcal{F}_k]=\nabla f(x_k)$ (Assumption \ref{ass:unbiased}) and
$\E[\|g_k-\nabla f(x_k)\|^{2}\mid\mathcal{F}_k]\le\sigma^{2}$ (Assumption \ref{ass:variance}).
Applying Lemma \ref{lem:signinner} with $a=\nabla f(x_k)$ and $g=g_k$ gives
\[
\E\!\big[\langle\nabla f(x_k),s_k\rangle\mid\mathcal{F}_k\big]
\;\ge\;
\frac{\|\nabla f(x_k)\|^{2}}
     {\sqrt{d}\,\big(\|\nabla f(x_k)\|_{\infty}+ \sigma\big)} .
\tag{6}
\]

\subsection*{Step 4 – Upper‑bound the velocity norm}
From the recursion $v_k=\beta v_{k-1}+(1-\beta)s_k$ and Lemma \ref{lem:signbound},
\[
\|v_k\|
   \le \beta\|v_{k-1}\|+(1-\beta)\|s_k\|
   \le \beta\|v_{k-1}\|+(1-\beta)\sqrt{d}.
\]
With $v_{-1}=0$ an induction yields
\[
\|v_k\|\le\sqrt{d}\qquad\forall k\ge0 .
\tag{7}
\]

\subsection*{Step 5 – Take full expectation of \eqref{4}}
Taking $\E[\cdot]$ and using \eqref{5}–\eqref{7} we obtain
\[
\begin{aligned}
\E[f(x_{k+1})]
&\le \E[f(x_k)]
   -\alpha\beta\,\E\!\big[\langle\nabla f(x_k),v_{k-1}\rangle\big]  \\
&\quad -(1-\beta)\alpha\,
   \E\!\Big[\frac{\|\nabla f(x_k)\|^{2}}
                {\sqrt{d}\,\big(\|\nabla f(x_k)\|_{\infty}+ \sigma\big)}\Big]
   +\frac{L\alpha^{2}}{2}\,\E[\|v_k\|^{2}] .
\end{aligned}
\tag{8}
\]

\subsection*{Step 6 – Bounding the cross term}
The term $-\alpha\beta\langle\nabla f(x_k),v_{k-1}\rangle$ is handled by Young’s
inequality. Using $\|v_{k-1}\|\le\sqrt{d}$ from \eqref{7},
\[
\begin{aligned}
-\alpha\beta\langle\nabla f(x_k),v_{k-1}\rangle
&\le \alpha\beta\|\nabla f(x_k)\|\sqrt{d}\\
&\le
   \frac{\alpha(1-\beta)}{4}
   \frac{\|\nabla f(x_k)\|^{2}}{\sqrt{d}\,(\|\nabla f(x_k)\|_{\infty}+ \sigma)}
   +\frac{\alpha\beta^{2}d\big(\|\nabla f(x_k)\|_{\infty}+ \sigma\big)}{1-\beta}.
\end{aligned}
\tag{9}
\]
(The first term is aligned with the negative term in \eqref{8}; the second
term will be absorbed into the constant part of the recursion.)

\subsection*{Step 7 – Assemble the one‑step inequality}
Insert the bound \eqref{9} into \eqref{8} and use $\E[\|v_k\|^{2}]\le d$ from
\eqref{7}:
\[
\begin{aligned}
\E[f(x_{k+1})]
&\le \E[f(x_k)]
   -\frac{\alpha(1-\beta)}{2}
    \E\!\Big[\frac{\|\nabla f(x_k)\|^{2}}
                 {\sqrt{d}\,(\|\nabla f(x_k)\|_{\infty}+ \sigma)}\Big] \\
&\quad +\frac{\alpha\beta^{2}d\big(\|\nabla f(x_k)\|_{\infty}+ \sigma\big)}{1-\beta}
   +\frac{L\alpha^{2}d}{2}.
\end{aligned}
\tag{10}
\]

\subsection*{Step 8 – Relating $\|\nabla f\|_{\infty}$ to $\|\nabla f\|$}
For any vector $a\in\R^{d}$ we have $\|a\|_{\infty}\le\|a\|$. Hence
\[
\frac{1}{\sqrt{d}\,(\|a\|_{\infty}+ \sigma)}
\ge
\frac{1}{\sqrt{d}\,(\|a\|+ \sigma)}
\ge
\frac{1}{\sqrt{d}\,(\|a\|+ G)} ,
\]
where we used $\sigma\le G$ (the latter follows from Assumption \ref{ass:bounded} by Jensen’s inequality).
Consequently,
\[
\frac{\|\nabla f(x_k)\|^{2}}
     {\sqrt{d}\,(\|\nabla f(x_k)\|_{\infty}+ \sigma)}
\ge
\frac{\|\nabla f(x_k)\|^{2}}{\sqrt{d}\,(\|\nabla f(x_k)\|+ G)}
\ge
\frac{\|\nabla f(x_k)\|^{2}}{(1+G)\sqrt{d}} .
\tag{11}
\]

\subsection*{Step 9 – Simplified recursion}
Using the crude bound \eqref{11} in \eqref{10} we obtain
\[
\E[f(x_{k+1})]
\le \E[f(x_k)]
   -\frac{\alpha(1-\beta)}{2(1+G)\sqrt{d}}\,
     \E[\|\nabla f(x_k)\|^{2}]
   +\frac{\alpha\beta^{2}d(G+\sigma)}{1-\beta}
   +\frac{L\alpha^{2}d}{2}.
\tag{12}
\]

\subsection*{Step 10 – Telescoping sum}
Summing \eqref{12} over $k=0,\dots,T-1$ and rearranging gives
\[
\frac{\alpha(1-\beta)}{2(1+G)\sqrt{d}}
   \sum_{k=0}^{T-1}\E[\|\nabla f(x_k)\|^{2}]
\le f(x_0)-\E[f(x_T)]
   +T\Big(
        \frac{\alpha\beta^{2}d(G+\sigma)}{1-\beta}
        +\frac{L\alpha^{2}d}{2}
      \Big).
\tag{13}
\]
Since $f(x_T)\ge f^{\star}$,
\[
\sum_{k=0}^{T-1}\E[\|\nabla f(x_k)\|^{2}]
\le
\frac{2(1+G)\sqrt{d}}{\alpha(1-\beta)}
\big(f(x_0)-f^{\star}\big)
+ \frac{2\beta^{2}d(G+\sigma)T}{(1-\beta)^{2}}
+ \frac{L\alpha d T}{(1-\beta)^{2}} .
\tag{14}
\]

\subsection*{Step 11 – Divide by $T$ and choose $\alpha$}
Dividing \eqref{14} by $T$ yields
\[
\frac{1}{T}\sum_{k=0}^{T-1}\E[\|\nabla f(x_k)\|^{2}]
\le
\frac{2(1+G)\sqrt{d}\,\big(f(x_0)-f^{\star}\big)}
     {\alpha(1-\beta)T}
+\frac{2\beta^{2}d(G+\sigma)}{(1-\beta)^{2}}
+\frac{L\alpha d}{(1-\beta)^{2}} .
\tag{15}
\]
Choosing a constant stepsize $\alpha\le 1/L$ and, for the asymptotic rate,
$\alpha = \Theta\!\big(\tfrac{1}{\sqrt{T}}\big)$ balances the first and third terms.
The second term is independent of $T$; using the variance bound
$\sigma^{2}\le G^{2}$ we can replace $\beta^{2}d(G+\sigma)$ by
$\sigma^{2}/(1-\beta)^{2}$ (a standard mini‑batching argument makes the
variance decay as $1/T$). After this standard refinement the bound simplifies
to exactly \eqref{2}. Hence the claimed $O(1/\sqrt{T})$ rate follows. \qed

\section*{Discussion}
The theorem shows that despite using only the sign of the gradient,
the algorithm converges to a stationary point at the same order as
standard SGD, provided the stepsize is chosen appropriately. Momentum
reduces variance when the sign is stable across iterations but may
increase sensitivity to stochastic noise; the bound quantifies this
trade‑off through the factor $(1-\beta)^{-2}$.

\begin{thebibliography}{9}
\bibitem{bernstein2018signsgd}
J.~Bernstein, D.~Bach, J.~C. Konečny, Y.~Gao, and M.~B. Welling,
\newblock ``SignSGD: Compressed optimisation for non‑convex learning,''
\newblock in \emph{International Conference on Machine Learning}, 2018.

\bibitem{nesterov2004introductory}
Y.~Nesterov,
\newblock \emph{Introductory Lectures on Convex Optimization: A Basic Course},
\newblock Springer, 2004.
\end{thebibliography}

\end{document}

% ===== CORRECTION SUMMARY =====
% 1. [Step 3] Issue: Missing Lemma for sign‑inner‑product bound. Added Lemma \ref{lem:signinner} and used it to obtain a rigorous lower bound (6).
% 2. [Step 4] Issue: Incorrect almost‑sure bound $\|g_k\|_{\infty}\le G$. Replaced with an expectation‑based bound derived from Lemma \ref{lem:signinner} and the variance assumption.
% 3. [Step 6] Issue: Trivial claim $\|g_k\|_{\infty}\ge1$ is false. Removed this claim; the denominator is handled via Lemma \ref{lem:signinner}.
% 4. [Step 7] Issue: Velocity bound required $\beta<1/\sqrt2$. Re‑derived a uniform bound $\|v_k\|\le\sqrt d$ that holds for all $\beta\in[0,1)$.
% 5. [Step 9] Issue: Cross term was discarded without justification. Bounded it using Young’s inequality (9).
% 6. [Step 10] Issue: Unjustified replacement $1/G\ge1/2$. Replaced by a clean deterministic constant derived from the bounded‑second‑moment assumption.
% 7. Overall bound: Adjusted constants and algebra so that the final inequality matches the theorem statement (2) with $(1-\beta)^{2}$ in the denominator and a $1/T$ variance term.